\section{Expected Value: What Is a Bet Worth Before It Resolves?}

\subsection{The Expected Monetary Value Calculation}

Expected value is the foundational arithmetic of decision-making under uncertainty, and one of the most dangerously misapplied concepts in entrepreneurship.

For a set of mutually exclusive outcomes $x_i$ occurring with probabilities $p_i$, the expected monetary value is:
\begin{equation}
\operatorname{EV}(X) = \sum_{i=1}^n p_i x_i
\end{equation}

Consider an educational founder evaluating whether to launch a new advanced cohort. Measured strictly in \textbf{incremental net cash relative to doing nothing}, the assessed outcomes are:

\begin{tcolorbox}[
  enhanced,
  colback=BayesSurface!60,
  colframe=BayesBorder,
  boxrule=0.5pt,
  arc=3pt,
  left=3pt,right=3pt,top=3pt,bottom=3pt,
  boxsep=1pt
]
\scriptsize
\begin{tabularx}{\linewidth}{X r r}
\toprule
\textbf{Scenario} & \textbf{Assessed Probability} & \textbf{Incremental Cash Payoff} \\
\midrule
Cohort succeeds & 30\% & +₹8 lakh \\
Cohort fails & 70\% & $-$₹2 lakh \\
\midrule
\multicolumn{2}{l}{\textbf{Expected Monetary Value}} & \textbf{+₹1 lakh} \\
\bottomrule
\end{tabularx}
\end{tcolorbox}

The calculation is straightforward:
\begin{equation*}
\operatorname{EV} = 0.30 \times (+\text{₹}8\text{ lakh}) + 0.70 \times (-\text{₹}2\text{ lakh}) = +\text{₹}1\text{ lakh}
\end{equation*}

Under these stated numbers, proceeding with the launch produces positive expected monetary value compared to doing nothing.

This arithmetic, however, is not an automatic instruction to launch. It is a strictly \textbf{conditional calculation}: \textit{if} the estimated probabilities and financial payoffs accurately model reality, and \textit{if} the founder can absorb the loss without impairing future choices, the mathematical average across thousands of hypothetical repetitions is positive.

In the physical world, the founder does not live across 1,000 parallel universes. If the loss of ₹2 lakh wipes out liquid family savings or exhausts runway before payroll, a positive-EV choice can be commercially suicidal.

\subsection{Expected Value Versus Expected Utility}

The marginal utility of capital is profoundly non-linear. Losing your final ₹2 lakh of cash reserves can cause personal bankruptcy, emotional devastation, and institutional death. Earning an additional ₹2 lakh when possessing ₹50 lakh in reserves creates negligible incremental change in your operational reality.

An economically rational actor maximizes expected utility rather than raw expected rupees:
\begin{equation}
a^* = \arg\max_{a \in \mathcal{A}} \mathbb{E}\big[U(W \mid a)\big]
\end{equation}
where $W$ is total wealth following action $a$, and $U(\cdot)$ represents the concave utility function of the founder.

Furthermore, authentic founders carry multi-attribute objective functions that include student learning outcomes, intellectual autonomy, reputational integrity, and mission alignment. Attempting to force these non-financial values into a single monetary spreadsheet produces deceptive models.

\textbf{Expected value allows us to rigorously compare potential financial trajectories; it does not define what a human being ought to care about.}

\subsection{Decision Trees and Asymmetric Payoffs}

Before committing non-trivial capital, diagram a complete decision tree:
\begin{enumerate}[leftmargin=12pt,itemsep=2pt,topsep=1pt,parsep=0pt]
  \item Which discrete strategic actions are physically available right now?
  \item Which external events can occur following each respective choice?
  \item What are the net cash and reputational consequences of each branch?
  \item What secondary operational decisions unlock once an event resolves?
  \item Which specific paths lead directly to absorbing barriers and ruin?
\end{enumerate}

The true value of entrepreneurial strategy lives in nodes 4 and 5. A small, tightly scoped pilot creates the valuable option to expand aggressively upon success. A rigid, 3-year commercial property lease eliminates the option to retreat upon failure.

\subsection{Preventing Double-Counting Fallacies}

When constructing payoff models, maintain strict accounting hygiene. If the ₹2 lakh downside already includes the founder's software development spend, do not subtract that spend a second time. If the upside payoff represents gross top-line revenue rather than net operating cash contribution, your model is deeply corrupted. Explicitly declare whether every node represents revenue, gross profit, net operating cash flow, or long-term equity value.
